\exercisesection{3}

\begin{enumerate}
\def\labelenumi{\arabic{enumi}.}
\tightlist
\item
  \begin{enumerate}
  \def\labelenumii{(\alph{enumii})}
  \tightlist
  \item
    设 A 是一个环（不必交换），\(I_{s}\) （相应地 \(I_{d}\) ）是 A 中没有左逆（相应地右逆）的元素组成的集合。证明下列条件等价：(1) \(I_{s}\) 中两个元素之和属于 \(I_{s}\) ；(2) \(I_{s}\) 是一个左理想；(3) 在由 \(\neq A\) 的诸左理想组成的集合中存在一个最大的左理想 J。若这些条件成立，则反环 \(A^{0}\) 也满足这些条件，\(I_{s} = I_{d} = J\) ，J 是 A 的唯一的极大（左或右）理想，因而是 A 的 Jacobson 根，每个元素 \(x \notin J\) 都可逆，且 A/J 是一个域（参看《代数》第 VIII 章 §6 第 3 节）。这时我们也说 A 是一个局部环。
  \end{enumerate}
\end{enumerate}

\begin{enumerate}
\def\labelenumi{(\alph{enumi})}
\setcounter{enumi}{1}
\tightlist
\item
  证明：在局部环中，除 0 与 1 外不存在其他幂等元。
\end{enumerate}

\begin{enumerate}
\def\labelenumi{\arabic{enumi}.}
\setcounter{enumi}{1}
\tightlist
\item
  \begin{enumerate}
  \def\labelenumii{(\alph{enumii})}
  \tightlist
  \item
    设 A 是一个（交换）局部环，其极大理想 m 是主理想，且满足 \(\bigcap_{n\geqslant1}m^{n}=0\) （参看第 III 章 §3 第 2 节命题 4 的推论）。证明 A 的既 \(\neq0\) 又 \(\neq A\) 的理想只有诸幂 \(m^{n}\) ；由此推出 A 是 Noether 环，并且*或者是一个离散赋值环（第 VI 章）*，或者是一个使 1 为不可分解幂等元的拟主理想环（《代数》第 VII 章 §1 习题 6）。
  \end{enumerate}
\end{enumerate}

\begin{enumerate}
\def\labelenumi{(\alph{enumi})}
\setcounter{enumi}{1}
\tightlist
\item
  设 A 是由在 0 的一个邻域中定义且连续、并在点 0 可微的实值函数在点 t = 0 处的芽组成的环（《一般拓扑学》第 I 章第 3 版 § 6 第 10 节）。证明 A 是一个局部环，其极大理想 m 由函数 \(j: t \mapsto t; m^{n} \neq 0\) （对所有 n）的芽生成，但 A 不是主理想整环。若 c 是函数 \(t \mapsto \exp(-1/t^{2})\) 的芽，p 是 A 的一个不含 c 的素理想，证明商环 B = A/p 是一个局部整环，它不是主理想整环，而其极大理想是主理想。
\end{enumerate}

¶ 3. 设 A 是一个局部环（不必交换，参看习题 1），m 是它的 Jacobson 根。

\begin{enumerate}
\def\labelenumi{(\alph{enumi})}
\item
  设 L 是一个自由左 A-模，\(x \neq 0\) 是 L 的一个元素。设 \((e_{\iota})_{\iota \in I}\) 是 L 的一个基，对于它，x 的 \(\neq 0\) 分量的个数尽可能地小；若 \(x = \sum_{\iota \in I} \xi_{\iota} e_{\iota}\) ，设 J 是 \(\iota \in I\) 中使 \(\xi_{\iota} \neq 0\) 的元素组成的集合。证明：诸 \(\xi_{\iota} (\iota \in J)\) 中任何一个都不属于 A 中由其余分量生成的右理想。
\item
  在 (a) 的假设与记号下，设 P、Q 是 L 的两个互补子模，且 \(x \in P\) ；对每个 \(\iota \in I\) ，写 \(e_{\iota} = y_{\iota} + z_{\iota}\) ，其中 \(y_{\iota} \in P\) ，\(z_{\iota} \in Q\) 。证明：诸 \(y_{\iota}\) （其中 \(\iota \in J\) ）与诸 \(e_{\iota}\) （其中 \(\iota \notin J\) ）构成的族是 L 的一个基。（利用 (a) 证明：若 \(z_{\kappa} = \sum_{\iota \in I} \zeta_{\kappa_{\iota}} e_{\iota}\) ，则分量 \(\zeta_{\kappa_{\iota}}\) （其中 \(\kappa \in J\) 且 \(\iota \in J\) ）必属于 m，并利用命题 4 的推论 2。）由此推出：存在 P 的一个自由子模，它包含 x 且是 P 的直因子。
\item
  证明每个投射 A-模 P 都是自由的（利用 Kaplansky 定理（《代数》第 II 章 § 2 习题 4），归结到 P 由可数多个元素生成的情形，并对该情形应用 (b)）。
\item
  给出一个局部环 A 和一个平坦 A-模 M 的例子，使 M 不是自由的且满足 mM = M（取 \(A = \mathbf{Z}_{(p)}\) ）。*（在第 III 章 § 3 中，我们将给出 Noether 局部环 A 上的非自由忠实平坦 A-模的例子。）*
\item
  设 A 是一个局部环，M 是一个有限生成平坦 A-模。证明 M 是自由 A-模（利用第 I 章 § 2 习题 23 (e)）。
\end{enumerate}

\begin{enumerate}
\def\labelenumi{\arabic{enumi}.}
\setcounter{enumi}{3}
\item
  设 A 是一个局部环（不必交换），m 是它的 Jacobson 根。设 M、N 是两个自由 A-模（不一定有限生成），\(u: M \rightarrow N\) 是一个同态，使得 \(1 \otimes u: M/mM \rightarrow N/mN\) 是双射。证明 u 是双射。（先证明 u 是满射，再证明 u 是单射，归结到 M 与 N 都是有限生成的情形，并分别利用命题 4 的推论 1 与命题 6。）
\item
  用例子说明：命题 7 不能推广到局部环 A 不是约化环的情形。
\item
  \begin{enumerate}
  \def\labelenumii{(\alph{enumii})}
  \tightlist
  \item
    设 M 是《代数》第 VII 章 §3 习题 22 (h) 中定义的 Z-模；设 T 是 M 的扭子模。证明：对每个素数 p，\(\mathrm{T}_{(p)}\) 作为 \(\mathrm{M}_{(p)}\) 的子模是 \(\mathrm{M}_{(p)}\) 的直因子，尽管 T 不是 M 的直因子。
  \end{enumerate}
\end{enumerate}

\begin{enumerate}
\def\labelenumi{(\alph{enumi})}
\setcounter{enumi}{1}
\tightlist
\item
  设 \(\mathbf{M}\) 是《代数》第 VII 章 § 3 习题 5 中定义的 \(\mathbf{Z}\) -模；证明：对每个素数 \(p\) ，\(\mathrm{M}_{(p)}\) 是自由 \(\mathbf{Z}_{(p)}\) -模，尽管 \(\mathbf{M}\) 不是自由 \(\mathbf{Z}\) -模。
\end{enumerate}

\begin{enumerate}
\def\labelenumi{\arabic{enumi}.}
\setcounter{enumi}{6}
\item
  设 A 是一个环，\((\mathrm{S}_{\lambda})_{\lambda \in L}\) 是 A 的乘性子集的一个族，使得对 A 的每个极大理想 m，存在 \(\lambda \in L\) 满足 \(m \cap S_{\lambda} = \varnothing\) （参看 §2 习题 8）。设 M、N 是两个 A-模，u: M → N 是一个同态。为使 u 是满射（相应地，单射、双射、零同态），必须且只需：对所有 \(\lambda \in L\) ，\(S_{\lambda}^{-1}u\) : \(S_{\lambda}^{-1}M \to S_{\lambda}^{-1}N\) 是满射（相应地，单射、双射、零同态）。第 3 节定理 1 的诸推论可以怎样推广？
\item
  设 A 是一个环，B 是一个 A-代数（不必交换），E 是一个有限生成左 B-模。
\end{enumerate}

\begin{enumerate}
\def\labelenumi{(\alph{enumi})}
\item
  设 S 是 A 的一个乘性子集，并给 \(S^{-1}B = S^{-1}A \otimes_{A} B\) 赋以它的 \(S^{-1}A\) -代数结构；于是可以把 \(S^{-1}E\) 视为一个左 \(S^{-1}B\) -模（《代数》第 VIII 章 § 7 第 1 节），它同构于 \(S^{-1}B \otimes_{B} E\) 。证明：若 E 是平坦 B-模，则 \(S^{-1}E\) 是平坦 \(S^{-1}B\) -模。
\item
  设 \((\mathbf{S}_{\lambda})_{\lambda \in L}\) 是 A 的乘性子集的一个族，使得对 A 的每个极大理想 m，存在 \(\lambda \in L\) 满足 \(m \cap S_{\lambda} = \varnothing\) 。证明 \(\bigoplus_{\lambda \in L} S_{\lambda}^{-1} B\) 是一个忠实平坦 B-模（左模或右模；参看 § 2 习题 8）。证明：若对所有 \(\lambda \in L\) ，\(S_{\lambda}^{-1} E\) 是平坦（相应地，忠实平坦）\(S_{\lambda}^{-1} B\) -模，则 E 是平坦（相应地，忠实平坦）B-模（利用第 I 章 § 3 习题 9，取 \(C_{\lambda} = S_{\lambda}^{-1} A\) ）。
\item
  假设诸 \(S_{\lambda}\) 满足条件 (b)，且 L 有限。证明：若对每个 \(\lambda \in L\) ，\(S_{\lambda}^{-1}E\) 是有限生成投射 \(S_{\lambda}^{-1}B\) -模，则 E 是有限生成投射 B-模。
\end{enumerate}

\begin{enumerate}
\def\labelenumi{\arabic{enumi}.}
\setcounter{enumi}{8}
\item
  证明：为使一个（交换）环 A 是绝对平坦的（第 I 章 § 2 习题 17），必须且只需对 A 的每个极大理想 m，\(A_{m}\) 都是域（注意绝对平坦局部环必为域，并利用命题 15 的推论）。
\item
  设 A、B 是两个环，\(\rho: A \rightarrow B\) 是一个同态，q 是 B 的一个素理想，\(p = \bar{\rho}^{-1}(q)\) 。假设存在一个 B-模 N，使得 \(N_{q}\) 是一个不化为 0 的平坦 A-模，且 \(N_{q}\) 是一个有限生成 \(A_{p}\) -模；则素理想 p 在 A 中是极小的。（注意 \(N_{p}\) 这时是一个忠实平坦 \(A_{p}\) -模，并利用 §2 第 5 节命题 11 的推论 4。）
\item
  设 A 是一个环，a 是 A 的一个理想，S 是由其在 A/a 中的典范像不是零因子的那些元素组成的 A 的乘性子集；设 \(\Phi\) 是 A 的与 S 不相交的理想组成的集合中的极大元组成的集合（因而诸理想 \(p \in \Phi\) 是素的，参看 §2 第 5 节命题 11）。证明：a 关于诸理想 \(p \in \Phi\) 的饱和的交就是 a 本身（归结到 a = 0 的情形，并利用定理 1 的推论 1）。
\end{enumerate}

¶ 12. 设 \(A = \prod_{i=1}^{n} A_i\) 是有限多个局部环 \(A_i\) 组成的族的乘积，这些局部环典范地等同于 \(A\) 的理想。设 \(B\) 是 \(A\) 的一个子环，使得 \(\text{pr}_i B = A_i\) 对 \(1 \leqslant i \leqslant n\) 成立。证明 \(B\) 是至多 \(n\) 个局部环的直复合，并且是 \(n\) 个局部环的直复合仅当 \(B = A\) 时才可能。（对 \(n\) 作归纳，考虑由 \(a_i = B \cap A_i\) 给出的 \(B\) 的理想。相继考察 \(a_i = A_i\) 至少对一个指标 \(i\) 成立的情形，以及 \(a_i \neq A_i\) 对所有 \(i\) 成立的情形。注意：若 \(a_i \neq A_i\) ，则 \(B\) 的每个极大理想都包含 \(a_i\) ，并且通过考虑环 \(B / a_i\) ，断言至多存在 \(n - 1\) 个属于 \(B\) 的互不相同的极大理想；若这时 \(B\) 不是局部环，则利用归纳假设与习题 1 (b) 归结到 \(a_j = A_j\) 对某个 \(j \neq i\) 成立的情形。）给出一个例子，其中 \(n > 1\) 且 \(B\) 是局部环，但不与任何 \(A_t\) 同构（取诸 \(A_t\) 为同一域 \(K\) 上的代数，其极大理想 \(m_t\) 的平方为零）。

¶ 13. 设 G 是一个有限群，其阶 n 是一个素数 p 的幂 \(p^{f}\) ，而 K 是特征为 p 的域。

\begin{enumerate}
\def\labelenumi{(\alph{enumi})}
\item
  设 E 是一个 G 在其上作用的有限集（《代数》第 I 章 § 7 第 2 节）；若 \(E'\) 是 E 中对所有 \(g \in G\) 都不变的元素组成的集合，证明 \(\text{Card}(E) \equiv \text{Card}(E') \pmod{p}\) 。由此推出：若 G 不化为它的单位元 e，则 G 的中心 Z 不化为 e（让 G 通过内自同构作用在它自身上）。断言 G 是可解的（《代数》第 I 章 § 6 习题 14）。
\item
  设 V 是 K 上的一个向量空间，\(\rho\) 是 G 到 \(\mathbf{GL}(V)\) 的一个同态。设 \(V'\) 是 \(x \in V\) 中在线性映射 \(\rho(g)\) （其中 \(g \in G\) ）下不变的元素组成的集合。证明：若 \(V'\) 化为 0，则 V 必化为 0。（若 \(x \in V\) 且 \(x \neq 0\) ，对由诸元素 \(\rho(g).x\) （其中 \(g \in G\) ）生成的 V 的加法子群 E 应用 (a)。）
\item
  设 \(\mathbf{A} = \mathbf{K}^{(\mathbf{G})}\) 是群 \(\mathbf{G}\) 在 \(\mathbf{K}\) 上的代数（《代数》第 III 章 § 1），\(\mathbf{I}\) 是 \(\mathbf{A}\) 的（ \(\mathbf{K}\) 上的）由诸元素 \(1 - g\) （其中 \(g \in \mathbf{G}\) ）生成的向量子空间。证明 \(\mathbf{I}\) 是 \(\mathbf{A}\) 的 Jacobson 根，且 \(\mathbf{A} / \mathbf{I}\) 同构于 \(\mathbf{K}\) （利用 (b) 及 Jacobson 根的定义）；由此推出 \(\mathbf{I}\) 是幂零的。
\item
  设 V 是一个有限生成 A-模，V′ 是 V 中对所有 \(g \in G\) 满足 g.x = x 的元素 x 组成的集合。证明下列不等式成立
\end{enumerate}

\[
(*) \dim_ {\mathbb {K}} (\mathrm{V}) \leqslant n. \dim_ {\mathbb {K}} (\mathrm{V} / \mathrm{IV}); \quad (* *) \dim_ {\mathbb {K}} (\mathrm{V}) \leqslant n. \dim_ {\mathbb {K}} (\mathrm{V} ^ {\prime}).
\]

（为建立 (*)，利用命题 4 的推论 2；通过考虑 V 在 K 上的对偶向量空间推出 (**)。）为使不等式 (*) （相应地 (**)）两边相等，必须且只需 V 是自由 A-模（参看命题 5 的推论 1）。

\begin{enumerate}
\def\labelenumi{\arabic{enumi}.}
\setcounter{enumi}{13}
\tightlist
\item
  设 A 是一个环，它不必交换，且其 Jacobson 根 m 使得 A/m 是一个域；设 M 是一个有限生成自由 A-模。
\end{enumerate}

\begin{enumerate}
\def\labelenumi{(\alph{enumi})}
\item
  设 \(\mathbf{M}'\) 是另一个有限生成自由 A-模，\(u: \mathbf{M} \to \mathbf{M}'\) 是一个同态，使得 \(u(\mathbf{M})\) 是 \(\mathbf{M}'\) 的一个直因子；数 \(\dim(\mathbf{M})\) 称为 \(u\) 的秩，记作 \(\operatorname{rg}(u)\) 。对这样的同态，\(\operatorname{Ker}(u)\) 是 \(\mathbf{M}\) 的直因子，且 \(\operatorname{rg}(u) = \dim(\mathbf{M}) - \dim(\operatorname{Ker}(u))\) ；为使 \(u\) 是单射（相应地满射），必须且只需 \(\operatorname{rg}(u) = \dim(\mathbf{M})\) （相应地 \(\operatorname{rg}(u) = \dim(\mathbf{M}')\) ）。转置 \(^t u\) 使得 \(^t u(\mathbf{M}'*)\) 是 \(\mathbf{M}^*\) 的直因子，且 \(\operatorname{rg}(^t u) = \operatorname{rg}(u)\) 。
\item
  设 \(n = \dim(\mathbf{M})\) 。\(\mathbf{M}\) 的维数为 1（相应地 \(n - 1\) ）的直因子称为直线（相应地超平面）。一个异于恒等映射的自同构 \(u \in \mathbf{GL}(\mathbf{M})\) 称为一个平延，如果存在超平面 \(\mathbf{H}\) 含于 \(\mathbf{M}\) ，它的所有元素在 \(u\) 下不变；这时 \(u(x) = x + a\phi(x)\) ，其中 \(\phi\) 是 \(\mathbf{M}\) 上的一个线性形式，使得 \(\mathbf{H} = \operatorname{Ker}(\phi)\) ，而 \(a \in \mathbf{H}\) ；证明其逆。若 \(\mathbf{A}\) 交换，证明每个行列式为 1 的自同构 \(u \in \mathbf{GL}(\mathbf{M})\) 都是平延之积（注意，在 \(u\) 关于 \(\mathbf{M}\) 的任意一组基的矩阵中，每一列至少含有 \(\mathbf{A}\) 的一个可逆元素，且形如 \(I + E_{ij} (i \neq j)\) 的矩阵是一个平延的矩阵）。
\item
  给出一个自同构 \(u \in \mathbf{GL}(\mathbf{M})\) 的例子，使 \(1 - u\) 的核不是 \(\mathbf{M}\) 的直因子（取 \(\mathbf{A}\) 为局部环 \(\mathbf{K}[[\mathbf{X}]]\) ，即域 \(\mathbf{K}\) 上一个未定元的形式幂级数环，并取 \(\mathbf{M} = \mathbf{A}\) ）。
\item
  给出 M 的直因子 N、P（必然是自由的）的例子，使得 \(N + P\) 与 \(N \cap P\) 都不是 M 的直因子。（取 A = K{[}X{]}/(X²)，其中 K 是一个域，M = A²。）
\end{enumerate}

¶ 15. 设 A 是一个（交换）环，M 是一个 A-模。

\begin{enumerate}
\def\labelenumi{(\alph{enumi})}
\item
  为使 M 的子模 \(M'\) 是纯的（第 I 章 §2 习题 24），必须且只需对 A 的每个极大理想 m，\(M_{m}'\) 都是 \(A_{m}\) -模 \(M_{m}\) 的纯子模（利用第 3 节定理 1）。
\item
  假设 A 是以 m 为极大理想的局部环，M 是一个有限生成自由 A-模。为使 M 的一个有限生成子模 \(M'\) 是纯的，必须且只需它是 M 的直因子。（利用第 2 节命题 5 的推论 1，归结为证明：若 \(M' \subset mM\) ，则 \(M'\) 仅当 \(M' = 0\) 时才可能是 M 的纯子模。）
\end{enumerate}

\begin{enumerate}
\def\labelenumi{\arabic{enumi}.}
\setcounter{enumi}{15}
\tightlist
\item
  设 \((\mathbf{A}_{\lambda}, f_{\mu \lambda})\) 是局部环的一个正向系，使得诸 \(f_{\mu \lambda}\) 都是局部同态；设 \(m_{\lambda}\) 是 \(\mathbf{A}_{\lambda}\) 的极大理想，\(\mathbf{K}_{\lambda} = \mathbf{A}_{\lambda} / m_{\lambda}\) 。则 \(\mathbf{A} = \varinjlim \mathbf{A}_{\lambda}\) 是一个局部环，其极大理想为 \(m = \varinjlim m_{\lambda}\) ，剩余域为 \(\mathbf{K} = \varinjlim \mathbf{K}_{\lambda}\) 。此外，若 \(m_{\mu} = \mathbf{A}_{\mu} m_{\lambda}\) 对 \(\lambda < \mu\) 成立，则 \(m = Am_{\lambda}\) 对所有 \(\lambda\) 成立。
\end{enumerate}

